% ═══════════════════════════════════════════════════════════════════
% Chapter 11 — Machine IR: IR003*
% ═══════════════════════════════════════════════════════════════════

\section{Machine IR: \IRstar{3}}
\label{sec:ir003}

\IRstar{3} is Crucible's vendor-machine-level intermediate representation: the form in which a compiled kernel is held between the per-vendor backend's lowering pass (Section~\ref{sec:mimic-template}) and the encoder that emits the vendor binary format.
\IRstar{3} is the deepest tier of Crucible's three-tier IR stack, sitting below \IR{2} (Chapter~\ref{sec:ir002}) and above the per-vendor binary bytes that the kernel driver consumes.

The asterisk in \IRstar{3} denotes a \emph{family} of representations rather than a single representation.
\IR{1} is one IR shared across the entire system; \IR{2} is one IR shared across the entire system; \IRstar{3} is six IRs --- one per accelerator vendor in the current design --- bound by a shared protocol but not by a shared type.
The set is open: adding a new vendor to Mimic adds one variant.

The reason for the split is structural.
At the operation level (\IR{1}) all vendors realize the same abstract semantics.
At the kernel level (\IR{2}) all vendors honor the same recipe, the same shape commitment, the same layout discipline.
At the machine level the abstractions cease to be vendor-uniform.
NVIDIA Hopper encodes 128-bit Mercury instruction words across sixteen format groups, with warpgroup-scoped \code{wgmma} on sm\_90 and asynchronous tensor-core MMA targeting tensor memory on sm\_100; AMD CDNA3 encodes a register-file-rich VLIW-like stream with wave-scoped accumulators driven by an in-order dispatcher and synchronized by \code{s\_waitcnt} vectors; Google TPU executables are protobuf-wrapped opaque payloads, reverse-engineered from the libtpu shipping binary, consumed by an on-chip scalar processor that schedules MXU, VPU, and DMA work; AWS Trainium NEFF binaries are tarballs of per-engine bytecode streams (TensorEngine, VectorEngine, ScalarEngine, GpSimd, DMA, Sync) with each engine's stream a statically scheduled sequence; Cerebras realizes spatial computation across a fabric of tiles described by Cerebras Software Language; the CPU backend writes native ELF directly or stages source for a JIT compiler.
There is no useful common type that admits all six simultaneously without becoming a thin wrapper over each.

The Crucible decision is therefore to commit to per-vendor IRs at this layer and to specify a uniform protocol that every vendor's IR satisfies.
The protocol is what makes Mimic's surface uniform; the IRs themselves are not.
This chapter specifies the protocol, the per-vendor variants, the recipe realization at the machine level, the encoder--decoder discipline, and the verification regime.
The misc/MIMIC.md design specification documents the per-vendor instruction encoders, the bit layouts, the EIATTR catalogs, the calibrated opcode tables, and the per-vendor lowering passes.

\subsection{Position in the IR stack}
\label{sec:ir003-position}

\IRstar{3} is the closest IR to silicon.
At the top of the stack, \IR{1} represents what the model computes.
\IR{2} narrows the representation to a kernel graph with shape, dtype, layout, and pinned recipe; \IR{2} deliberately commits no vendor-specific tile shapes, no vendor address spaces, no register allocation, and no instruction schedule (Section~\ref{sec:ir002-not}).
\IRstar{3} commits all of those.

Specifically, every \IRstar{3} program contains:

\begin{itemize}
  \item A linear sequence of vendor-specific machine instructions in canonical SSA form (\emph{CSSA}, Section~\ref{sec:ir003-cssa}), each carrying an opcode, operands, an output destination, scheduling-slot metadata, scoreboard or wait-counter annotations, and any vendor-specific decoration the ISA requires.
  \item Address-space tags on every memory operand identifying which tier of the per-vendor memory hierarchy the load or store touches (registers, fast-local-storage, distributed-shared, tensor memory, main memory, constant bank, instruction cache, special-purpose buffers).
  \item A complete register allocation: the mapping from logical SSA values to physical registers including spill placements, bank-conflict avoidance, and the live-range layout that determines occupancy.
  \item An instruction schedule: cycle-or-stage-accurate ordering, pipe assignments, scoreboard-slot reservations, dual-issue groupings, and the placement of synchronization primitives (mbarrier waits, named-barrier arrivals, tensor-memory commit groups, wait-counter sequences, fabric-routing semaphores, and so on).
  \item Memory-fence and synchronization primitives required by the kernel's data flow and by the recipe's determinism tier.
  \item All metadata required by the binary format: register counts, fast-local-storage usage, occupancy hints, debug-info entries, ELF-attribute tags, relocation records.
\end{itemize}

\IRstar{3} therefore is the form of the program that maps one-to-one onto bytes in the per-vendor binary format (cubin on NVIDIA, HSACO on AMD, executable on TPU, NEFF on Trainium, CSL on Cerebras, native ELF on CPU).
Encoding \IRstar{3} produces those bytes; decoding the bytes reconstructs \IRstar{3}.
The next step beyond \IRstar{3} is hardware execution; there is no further IR.

\subsection{Common discipline}
\label{sec:ir003-common}

Although the variants differ per vendor, they share a uniform discipline that the Mimic shared core can rely on without knowing the vendor-specific details.
The discipline is enumerated in misc/MIMIC.md \S2 as a set of contractual obligations on every backend; the architecture-level summary is four properties.

\begin{property}[Canonical SSA form]
\label{prop:ir003-cssa}
Every \IRstar{3} program is in CSSA: every value has exactly one definition, every instruction reads at most a fixed bounded number of operands per the vendor's ISA, and the program contains no unstructured control flow other than what the vendor's ISA models natively (predication, branch-divergence-aware control, partial-warp predicates, and similar).
The discipline allows the per-vendor register allocator and instruction scheduler to operate over a uniform abstract substrate even when the concrete instruction set differs.
\end{property}

\begin{property}[Address-space tagging]
\label{prop:ir003-as-tags}
Every memory operand carries the address-space tag identifying the storage tier it accesses.
The set of tags is per-vendor (registers, shared memory, distributed shared memory, tensor memory, global memory, constant memory, instruction cache on NVIDIA; the analog set per backend).
Crucible exposes a uniform abstract address-space taxonomy at the \IR{2}-to-\IRstar{3} boundary; per-vendor lowering refines abstract tags into concrete vendor tags.
\end{property}

\begin{property}[Scheduling-slot annotation]
\label{prop:ir003-sched}
Every instruction carries the scheduling metadata the vendor's hardware requires.
On NVIDIA Hopper, this is the per-instruction stall count from the SASS control field, the scoreboard wait mask, the read-scoreboard slot, the write-scoreboard slot, and the yield and reuse hint bits.
On AMD CDNA3, this is the wait-counter vector (vmcnt, lgkmcnt, expcnt) and the wave-pipeline placement.
On TPU, this is the explicit cycle index in the static schedule and the pipe assignment.
The scheduling metadata is structurally part of \IRstar{3} and survives encoding to the binary format.
\end{property}

\begin{property}[Round-trip-correct binary encoding]
\label{prop:ir003-roundtrip}
For every valid \IRstar{3} program $p$ on backend $V$,
\[
    \mathrm{decode}_V(\mathrm{encode}_V(p)) \;=\; p
\]
holds modulo a documented set of normalization rewrites that the per-vendor decoder applies (re-canonicalization of equivalent encodings of the same instruction, alignment-padding nops elision, debug-only metadata that the binary does not preserve).
The property is asserted continuously by the per-vendor encoder--decoder regression suite.
A pull request that ships an encoder change without the corresponding decoder change, or vice versa, fails the suite and the build is rejected.
\end{property}

The CSSA discipline is what makes the shared simulator framework, the MAP-Elites driver, and the insight extractor (Chapter~\ref{sec:mimic}) operate uniformly across backends.
The address-space tagging is what makes the cross-vendor numerics CI matrix (Chapter~\ref{sec:cross-vendor-ci}) able to compile the same \IR{2} program through every backend and verify that each backend's recipe realization stays within the recipe's declared tolerance.
The scheduling-slot annotation is what makes \IRstar{3} a faithful proxy for the binary's runtime behavior, and therefore what makes the simulator's medium and accurate tiers calibratable.
The round-trip property is what allows the L2 KernelCache (Chapter~\ref{sec:cipher}) to load a cached binary back into a usable \IRstar{3} representation for Augur's prediction path.

\subsection{Canonical SSA form}
\label{sec:ir003-cssa}

The IR is in single-static-assignment form by construction; CSSA additionally requires five structural properties.

\begin{enumerate}
  \item \textbf{Definitions are unique by name.}
        Every SSA value is identified by a per-vendor strong-typed identifier (a register-name strong type on NVIDIA, the analog on each other backend).
        The identifier is arena-bounded; no two definitions in the same program share an identifier.

  \item \textbf{Operand reads are explicit.}
        Every instruction declares its full set of input operands as named-value references.
        There are no implicit reads (the NVIDIA status registers, the AMD VCC and EXEC masks, the TPU prefetch register file) other than those modeled as named values produced by upstream instructions.

  \item \textbf{Predication is a value, not a flag.}
        Predicated execution (per-thread predicate registers on NVIDIA, the EXEC mask on AMD, per-lane predicates on TPU) is modeled as an SSA value produced by a predicate-defining instruction and consumed as a typed operand by the predicated instructions.

  \item \textbf{Side effects are explicit.}
        Memory writes, shared-memory writes, fence operations, named-barrier arrivals, and tensor-memory commits are SSA-formed: each carries an explicit ordering token consumed by the subsequent fence or wait.
        No backend uses a global ``side-effect frame'' or implicit ordering.

  \item \textbf{Scope structure follows the vendor's hardware model.}
        On NVIDIA the relevant scopes are thread, warp, warpgroup, CTA, cluster (CGA), and grid; on AMD lane, wave, workgroup, and grid; on TPU lane, sublane, tile, and grid; on Trainium per-engine plus device plus fleet; on Cerebras fabric-tile, fabric-region, and wafer.
        Each \IRstar{3} variant exposes the scopes its ISA exposes.
\end{enumerate}

Property~\ref{prop:ir003-cssa} is the structural commitment.
Per-vendor variants additionally publish a verifier (the per-vendor analog of \code{Ir003NvVerifier}) that the Mimic shared core invokes after every lowering pass and after every MAP-Elites mutation that touches the program structurally.
A program that fails verification cannot be encoded; the search loop discards it and continues.

\subsection{The per-vendor variants}
\label{sec:ir003-variants}

The current set of \IRstar{3} variants is six.
Table~\ref{tab:ir003-variants} summarizes the variants at a glance; the per-vendor discussion that follows elaborates each.

\begin{table}[h]
\centering
\footnotesize
\begin{tabular}{l l l l l}
\toprule
\textbf{Variant} & \textbf{Vendor / family} & \textbf{Tensor-core family} & \textbf{Binary format} & \textbf{Driver protocol} \\
\midrule
\textsc{ir003nv}  & NVIDIA Hopper, Blackwell  & WGMMA, tcgen05 mma   & cubin (ELF + EIATTR)         & \code{/dev/nvidia*} ioctl + UVM \\
\textsc{ir003am}  & AMD CDNA3+, RDNA3+        & MFMA on \code{ACC\_VGPR}   & HSACO (AMDGPU ELF)           & \code{/dev/kfd} + amdgpu KMD \\
\textsc{ir003tpu} & Google TPU v5+, v6e, v7   & MXU systolic array   & TPU executable (proto + payload) & \code{/dev/accel*} ioctl plane \\
\textsc{ir003trn} & AWS Trainium 1, 2, 3      & TensorEngine + GpSimd & NEFF (per-engine tarball)         & \code{/dev/neuron*} ioctl plane \\
\textsc{ir003cer} & Cerebras WSE3+            & Wafer fabric routing & CSL (on-fabric program)     & \code{/dev/cerebras*} plane \\
\textsc{ir003cpu} & x86\_64, AArch64, RISC-V  & Scalar FMA + AVX-512 / NEON / SVE & native ELF \code{.so} or JIT & direct execution \\
\bottomrule
\end{tabular}
\caption{Per-vendor \IRstar{3} variants. Each variant is the full machine IR for one accelerator family and is owned by the corresponding Mimic backend.}
\label{tab:ir003-variants}
\end{table}

\subsubsection{IR003NV --- the reference-complete variant}
\label{sec:ir003-nv}

The most-elaborated variant in the current implementation is \textsc{ir003nv}.
It encodes the Hopper and Blackwell SASS family as a CSSA-form sequence over an aggressively-normalized opcode set.
Each instruction carries: an opcode encoded as packed major / minor / sub-op fields, up to three input operands, up to two output destinations (a tensor-core MMA has both a destination and an accumulator), a 4-bit predicate, a stall count drawn from the SASS control field, a 6-bit scoreboard wait mask, a 3-bit read-scoreboard slot, a 3-bit write-scoreboard slot, yield and reuse hint bits, an asynchronous-issue flag, and a parametric-shape flag for symbolic tile dimensions.
Address-space tags on memory operands distinguish registers, shared memory, distributed shared memory, tensor memory, global memory, and constant bank.

The encoder produces 128-bit Mercury instruction words organized into the sixteen format groups documented for sm\_75 and later, packages them into a cubin ELF, and emits the EIATTR catalog the NVIDIA driver requires for \code{cuModuleLoad}.
The EIATTR catalog includes register counts, frame and stack sizes, kernel-parameter constant-bank metadata, launch-bound annotations, named-barrier and mbarrier counts, cluster-launch geometry, the tcgen05 1-CTA-or-2-CTA-used flag (sm\_100+), the sparse-MMA mask, exit-instruction offsets, the toolkit version stamp, and the Mercury ISA version stamp; the exhaustive per-function value derivation is documented in misc/MIMIC.md \S15.

The decoder produces a \struct{DecodedProgram}, a simulator-facing form arena-owned by Mimic.
\struct{DecodedProgram} deliberately strips Forge-side metadata (pre-lowering attributes, mutation-source tracking, scheduling-group labels) that does not round-trip through the binary; the simulator does not need it.
The encoder--decoder pair is derived from a single per-instruction encoding table, ensuring that the two cannot disagree about the bit layout of any instruction.

The simulator subsystem in the NVIDIA backend specializes the shared framework with the GTO warp scheduler, the four-level memory hierarchy (registers, shared memory, L2, HBM), the WGMMA and TMA async pipelines on Hopper, and the tcgen05 with TMEM async pipeline on Blackwell.
The runtime library wraps \code{/dev/nvidia*}, \code{/dev/nvidia-uvm}, and \code{/dev/nvidiactl} ioctls, replacing libcuda, libcudart, libnvJitLink, and libnvrtc.
The collective library implements ring, tree, halving-doubling, and hierarchical algorithms over NVLink intra-node and RDMA inter-node, replacing NCCL and the upper layers of UCX.

\subsubsection{IR003AM --- AMD CDNA and RDNA}
\label{sec:ir003-am}

\textsc{ir003am} encodes the AMDGPU ISA for CDNA3 (gfx942 on MI300X, gfx950 on MI325X) and CDNA4, plus RDNA3 and later for consumer parts.
The structural pattern matches \textsc{ir003nv}: CSSA form, per-instruction wait-counter vector (\code{vmcnt}, \code{lgkmcnt}, \code{expcnt}), VGPR / SGPR / AGPR allocation, LDS tags on memory operands, MFMA pipeline placement on \code{ACC\_VGPR} accumulator registers.
The wait-counter vector is the AMD-specific scheduling discipline: an \code{s\_waitcnt} instruction stalls the wave until the counters drop below specified thresholds, replacing the per-instruction scoreboard slots that NVIDIA uses.

The encoder produces AMDGPU machine code organized into the HSACO container format (AMDGPU ELF with rodata kernel descriptors and DWARF debug-info sections).
The encoder may reuse LLVM-AMDGPU internally for ISA-table selection (the generator tables are open and stable); the HSACO wrapper, the scheduler, and the peephole pass are Mimic's own.
The decoder is the inverse.

The runtime library wraps \code{/dev/kfd} (the kernel fusion driver shared by all AMDGPU devices on a host) and per-device \code{/dev/dri/renderD*} nodes.
The collective library targets the XGMI fabric intra-node and RoCE inter-node, replacing RCCL and the upper layers of \code{rocshmem}.
The amdgpu kernel module is the only AMD-vendor-proprietary component in the runtime path.

\subsubsection{IR003TPU and IR003TRN --- statically scheduled accelerators}
\label{sec:ir003-tpu-trn}

The TPU and Trainium variants differ structurally from the GPU variants in two ways.

\textbf{Static scheduling.}
Both architectures use cycle-accurate static scheduling: the IR commits to a precise cycle-by-cycle (or stage-by-stage) order at compile time rather than relying on a hardware scoreboard at runtime.
The CSSA form therefore carries explicit cycle indices and pipe assignments, and the simulator's medium tier becomes effectively cycle-accurate by construction (the schedule \emph{is} the simulation).
The compensating discipline is that the lowering pass must commit to a schedule that respects all dependency latencies; a schedule that is too aggressive produces a binary that hangs the device.

\textbf{Tile-granularity compute.}
On both architectures the unit of work is a tile-sized batched matrix-multiply or convolution, not an individual scalar instruction.
The IR's grain is therefore coarser, and the encoder produces tile-program bytecode rather than scalar instruction streams.

For TPU, the encoder produces a TPU executable: a protobuf-wrapped header plus an opaque payload consumed by the on-chip scalar processor.
The scalar processor's instruction set is reverse-engineered from libtpu (Mimic does not link against libtpu); the opcode families include MXU dispatch, VPU dispatch, VMEM load, DMA issue, sync barrier, and loop primitives.
Submission to the device is via DMA descriptor chain rather than a host-engine pushbuffer; the TPU autonomously walks the descriptor chain.

For Trainium, the encoder produces NEFF (Neuron Executable File Format): a tarball with per-engine bytecode streams.
The six engines are TensorEngine (the systolic-array tensor core), VectorEngine (per-lane SIMD), ScalarEngine (control flow, address arithmetic, scalar reductions), GpSimd (general-purpose SIMD for ops outside the tensor-core family), DMA (HBM and SBUF transfer descriptors), and Sync (cross-engine barrier coordination).
Submission is via \code{NEURON\_IOCTL\_SUBMIT\_BATCH} with the per-engine streams enqueued together; the on-chip scheduler coordinates the engines through Sync's barrier primitives.

Recipe realization on these statically scheduled variants follows the same discipline as on the GPU variants: tensor-core fragments are sized per the recipe's declared determinism tier, accumulator dtype is pinned, FMA contraction is configured per the recipe, and the cross-vendor numerics CI matrix verifies the realization.

\subsubsection{IR003CER --- spatial computation}
\label{sec:ir003-cer}

The Cerebras variant differs from all others in modeling \emph{spatial} rather than temporal computation.
The IR commits to fabric-tile placement (which on-chip tile owns which weight shard), on-fabric routing (which fabric paths carry which intermediate values), and per-tile CSL kernels rather than to a per-instruction schedule on a single chip.
The structural template is preserved (CSSA form, address-space tagging, round-trip-correct encoding), but the scope hierarchy is fabric-tile, fabric-region, and wafer rather than thread, warp, and CTA.

The encoder produces Cerebras Software Language binary format together with the on-fabric program.
The runtime library wraps \code{/dev/cerebras*}.
The collective library targets the SwarmX wafer fabric, with the in-fabric aggregation surface exposed as a NetworkOffload provider in CNTP (Chapter~\ref{sec:cntp}).

Cerebras is the most-different variant from the GPU baseline; its presence in Mimic demonstrates that the per-vendor IR template tolerates substantial architectural variation.

\subsubsection{IR003CPU --- the reference oracle}
\label{sec:ir003-cpu}

The CPU variant is the simplest because the binary target is the host's native ELF and the runtime path is direct execution.
The encoder writes either a native ELF shared object (loaded via \code{dlopen}) or stages C++ source plus invokes a JIT compiler; both modes are supported.
Address-space tags are minimal (registers, cache, DRAM); the schedule is whatever the host's out-of-order pipeline produces from a clean instruction stream.
The runtime library uses standard \code{malloc}, \code{mmap}, and pthreads; the collective library uses shared memory intra-node and CNTP TCP transport inter-node.

The CPU variant's role is the load-bearing one: it is the \emph{reference oracle} for the cross-vendor numerics CI matrix.
Its discipline is therefore inverted relative to the GPU and accelerator backends.
It commits to one specific scalar-FMA semantic: IEEE-754 with round-to-nearest-even, no contraction other than the recipe-declared FMA, no flush-to-zero unless the recipe requires it.
It produces, under \code{BITEXACT\_STRICT}, exactly the canonical pairwise-tree reduction sorted by index that every recipe-tier requires.
It makes no attempt at vectorization or instruction-scheduling tricks that would jeopardize bit-equivalence under \code{BITEXACT\_STRICT}.
The CPU variant must be available unconditionally on every developer machine and every CI runner; the cross-vendor matrix's pairwise comparison anchors on the CPU output.

\subsection{Recipe realization at the \IRstar{3} level}
\label{sec:ir003-recipe-realization}

The \struct{NumericalRecipe} attached to every \struct{KernelNode} (Section~\ref{sec:ir002-recipe}) is realized in \IRstar{3}.
Each per-vendor backend exposes a \code{realize\_recipe} entry point that, given a recipe and the per-vendor \struct{TargetCaps}, produces the concrete sequence of \IRstar{3} instructions implementing the recipe.
The realization is per-vendor, deterministic, and verifiable; the structural commitments are four.

\begin{enumerate}
  \item \textbf{Reduction structure.}
        The recipe declares one of \code{UNORDERED}, \code{ORDERED}, \code{BITEXACT\_TC}, or \code{BITEXACT\_STRICT} (Chapter~\ref{sec:numerical-recipe}).
        Realization emits the matching reduction structure: free-form (\code{UNORDERED}); per-warp tree-reduction with pinned topology (\code{ORDERED}); the recipe-declared tensor-core fragment with K bounded by the cross-vendor-equivalence threshold plus a pinned scalar outer reduction (\code{BITEXACT\_TC}); a canonical pairwise binary tree of scalar FMA operations sorted by lane index (\code{BITEXACT\_STRICT}).
  \item \textbf{Accumulator dtype.}
        The recipe pins the accumulator's representation; realization emits the matching tensor-core, MFMA, or MXU instruction with that accumulator dtype.
  \item \textbf{FMA contraction.}
        If the recipe forbids FMA contraction across statements (the IEEE-754 default, used under \code{BITEXACT\_STRICT}), realization breaks fused-multiply-add patterns into separate multiply and add instructions.
        If the recipe allows contraction within a statement, realization emits FMA where the per-vendor scheduler benefits from it.
  \item \textbf{Rounding mode and denormal handling.}
        The recipe declares the rounding mode and the flush-to-zero policy; realization pins them via the per-vendor mode register (\code{setreg .denormal\_mode} on NVIDIA, \code{s\_setreg MODE} on AMD, the analog on each other backend).
\end{enumerate}

The per-vendor realization of \code{BITEXACT\_TC} deserves explicit attention because it is the recipe tier that delivers cross-vendor reproducibility at near-tensor-core throughput.
On NVIDIA, the realization prefers \code{wgmma m64n128k8} (or \code{m64n64k8}, \code{m64n32k8}) over the deeper \code{m64n128k16}, \code{m64n128k32}, and \code{m64n128k64} variants, because the carry-save-adder topology of summation trees beyond K=8 is vendor-specific and does not match AMD, TPU, or Trainium.
On AMD, the realization prefers \code{MFMA\_32x32x8\_F16} over the deeper K=16 variant for the same reason.
On TPU, the realization chunks large-K reductions into K=8 blocks with explicit accumulator resets between blocks.
On Trainium, the realization chunks similarly using PSUM accumulator resets at the K=8 cadence.
On all backends, the outer reduction across the K-fragments is emitted as a scalar FMA chain in pinned pairwise order; the FTZ flag is pinned via the mode register; the block-issue order is deterministic via the per-vendor launch-geometry pinning.

The cross-vendor numerics CI matrix (Chapter~\ref{sec:cross-vendor-ci}) verifies the realization.
For every (\struct{KernelKind}, \struct{NumericalRecipe}, target) triple, the matrix compiles the kernel through every per-vendor backend, executes it on real silicon (or on the simulator when hardware is unavailable), and compares the result pairwise against the CPU reference oracle.
A backend whose realization produces a result outside the recipe's declared tolerance fails the build.
\IRstar{3} is therefore the layer at which the cross-vendor portability contract is mechanically discharged: \IR{2} declares the recipe, \IRstar{3} realizes it, and the CI matrix verifies that the realization is faithful.

\subsection{The encoder--decoder pair}
\label{sec:ir003-codec}

Each per-vendor backend ships a paired encoder and decoder.
The encoder maps an \IRstar{3} program to bytes in the vendor's binary format; the decoder maps bytes to an \IRstar{3} program.
The two are derived from a single per-vendor instruction-encoding table, ensuring that the encoder and decoder cannot disagree about the bit layout of any instruction.

The encoder is the principal hot path of Forge's Phase~J (Section~\ref{sec:forge-phase-j}): every kernel produced by MAP-Elites passes through the encoder before commit.
The encoder is structured as a small set of primitives (the bitfield packer, the format-group dispatcher, the operand encoders, the control-field writer) plus larger components for the binary-format wrapping (cubin ELF emission, EIATTR catalog generation, DWARF line-table emission).
The bit-level details are documented in misc/MIMIC.md \S15.

The decoder's primary role is not loading kernels back into Crucible (Mimic's runtime library loads compiled kernels by content hash through a binary loader that does not need the IR back); it is enabling the simulator and Augur's prediction path to reason about cached or third-party-supplied binaries.
A binary loaded from the L3 KernelCache (Chapter~\ref{sec:cipher}) is decoded back to a simulator-facing form (\struct{DecodedProgram}) so the simulator can predict its performance under a new \struct{TargetCaps} during cross-chip migration.

The round-trip property (Property~\ref{prop:ir003-roundtrip}) is asserted continuously by a per-vendor regression suite that generates random valid \IRstar{3} programs through a per-vendor seeded fuzzer, encodes and decodes each, and compares structurally.
The suite is part of every PR build; encoder--decoder agreement is non-negotiable.

\subsection{Verification at the \IRstar{3} level}
\label{sec:ir003-verification}

In addition to the per-vendor CSSA verifier (Section~\ref{sec:ir003-cssa}), each backend implements four further verifiers.

\textbf{Register-pressure verifier.}
Asserts that the program does not exceed the chip's per-thread register file capacity nor force the program below the recipe's declared minimum occupancy.
A program that violates the constraint is rejected before encoding; MAP-Elites discards the candidate and continues.

\textbf{Scheduling verifier.}
Asserts every scoreboard or wait-counter read has a matching upstream write within the vendor's stall window, every named-barrier arrival has a matching wait, and every async-pipeline commit has a matching wait downstream.
A scheduling violation produces a binary that hangs the device; rejection at compile time is the structural mechanism that prevents this from reaching the runtime.

\textbf{Recipe-realization verifier.}
Given the kernel's pinned recipe, asserts the emitted instruction sequence implements the declared reduction structure, accumulator dtype, FMA-contraction policy, and rounding mode.
A program that fails recipe-realization verification is rejected before encoding; the cross-vendor numerics CI matrix is therefore guarded by a structural check at the per-program level in addition to the cross-vendor pairwise comparison.

\textbf{Native-recipe-support verifier.}
Asserts that the kernel's recipe is on the per-vendor backend's \code{native\_on} bitmap from the recipe registry (Chapter~\ref{sec:numerical-recipe}).
If the recipe is not natively supported on the target, the build is rejected at the \IR{2}-to-\IRstar{3} boundary before lowering proceeds; Forge's recipe selector should have declined the recipe at Phase~E, and the boundary check is a defense in depth.

These verifiers run in Forge's Phase~H (Section~\ref{sec:forge-phase-h}) immediately after per-vendor lowering and again after every MAP-Elites mutation that touches the program structurally.
The cost is approximately one to two milliseconds per kernel and is part of the per-kernel compilation budget.

\subsection{Properties}
\label{sec:ir003-properties}

The properties below are commitments of the \IRstar{3} specification, applying uniformly across all per-vendor variants.

\begin{property}[Per-vendor specialization with uniform protocol]
\label{prop:ir003-protocol}
Each \IRstar{3} variant is per-vendor and need not share a type with any other variant.
However, each variant satisfies the same protocol obligations: CSSA form, address-space tagging, scheduling-slot annotation, round-trip-correct encoding, recipe realization, and the verification suite.
A new vendor onboards Mimic by implementing the protocol; no other backend or shared-core code changes.
\end{property}

\begin{property}[Faithful proxy for the binary]
\label{prop:ir003-faithful}
For every per-vendor backend, the bit layout of the binary is determined by \IRstar{3}: encoding is a function of the \IRstar{3} program plus the binary-format header constants required by the loader.
Two \IRstar{3} programs that differ in any structural respect produce different binaries; two \IRstar{3} programs that differ only in fields the binary format does not preserve produce identical binaries.
The encoder is therefore a complete specification of the binary format from the IR's perspective.
\end{property}

\begin{property}[Recipe realization is verifiable per program]
\label{prop:ir003-recipe-verifiable}
For every \IRstar{3} program $p$ on backend $V$ targeting recipe $r$, the realization is verifiable in time linear in the size of $p$ by inspecting the emitted instruction stream against $r$'s declared structural properties.
A program that fails recipe-realization verification at the \IRstar{3} level is rejected before encoding, ensuring that no compiled binary can violate its declared recipe contract.
\end{property}

\begin{property}[CPU reference oracle is bit-exact]
\label{prop:ir003-cpu-bitexact}
The \textsc{ir003cpu} variant produces, under \code{BITEXACT\_STRICT}, scalar-FMA results with canonical pairwise-binary-tree reduction sorted by index, IEEE-754 round-to-nearest-even, no contraction, and no flush-to-zero.
The result is platform-independent (x86\_64, AArch64, and RISC-V produce identical bytes) by construction of the CPU lowering pass.
The property is the ground truth that the cross-vendor CI matrix verifies every other backend's realization against.
\end{property}

\begin{property}[Decoder reconstructs simulator-facing form]
\label{prop:ir003-decode-faithful}
For every per-vendor binary $b$ produced by the per-vendor encoder, the per-vendor decoder produces a simulator-facing form sufficient to drive every simulator tier (fast, medium, accurate).
The decoder may discard Forge-side metadata not needed for simulation, but it preserves all per-instruction structural information the simulator reads (opcode, operands, scheduling-slot, scoreboard or wait-counter, address-space, predicate).
\end{property}

\subsection{What \IRstar{3} does not contain}
\label{sec:ir003-not}

\IRstar{3} is the deepest IR but it is still an IR, not the silicon.
The boundary is structural and explicit: \IRstar{3} commits the static, deterministic, observable-from-the-IR properties of the program; it does not commit the runtime-contended, microarchitectural-state-dependent, or below-software properties.

\begin{itemize}
  \item \textbf{Transient hardware state.}
        Branch-prediction history, cache-line replacement state, fetch-pipeline pressure, and similar microarchitectural dynamics are not modeled in \IRstar{3}; they are properties of the runtime execution that the simulator approximates statistically.

  \item \textbf{Dynamic dispatch decisions.}
        On chips with dynamic warp or wave scheduling (NVIDIA's GTO, AMD's wave dispatch), the actual order in which warps progress at runtime is not pinned by \IRstar{3}; the schedule \IRstar{3} commits is the structural one (which instructions are dispatchable at which scoreboard or wait-counter state), not the temporal one (which warp gets the issue slot in cycle $t$).

  \item \textbf{Memory-system contention outcomes.}
        L2 cache evictions, HBM bank conflicts, fabric arbitration, and refresh-induced stalls are runtime properties; \IRstar{3} commits the address-space tags and access patterns but not the contention outcomes.

  \item \textbf{Below-software effects.}
        Temperature-dependent DRAM timing, voltage droop under transients, process variation across chips of the same SKU, ECC scrub overhead under non-microbenchmark workloads --- all transistor-level effects invisible at the IR level.
\end{itemize}

These omissions are why the simulator's medium and accurate tiers exist: \IRstar{3} pins enough to compute deterministic-execution properties exactly (cycle counts of the dependency-graph critical path, register-pressure timeline, scoreboard-stall structure) and enough to bound runtime-contended properties (HBM bandwidth ceiling, L2 hit-rate envelope), but it does not pretend to predict exact runtime behavior.
That is the simulator's job, calibrated against measurement, with drift attributed by Augur and corrected by Mimic's recalibration entry point.

\subsection{Relationship to the rest of the IR stack}
\label{sec:ir003-stack-summary}

The three IRs together form Crucible's compilation pipeline.
\IR{1} represents what the model computes.
\IR{2} represents how the computation is organized into kernels with shape, dtype, layout, and recipe commitments, in a form that any vendor could lower.
\IRstar{3} represents the per-vendor machine program that maps onto the silicon.
Each IR is content-addressable, each is independently cacheable, and each can be inspected, verified, and mutated.

The L1, L2, and L3 KernelCache hierarchy (Chapter~\ref{sec:cipher}) keys on the three IRs respectively: L1 keys on \IR{2} content hashes (federation-shareable across vendor mixes), L2 keys on \IRstar{3} content hashes (cross-chip-shareable within a vendor family), and L3 keys on the per-chip binary hash (chip-specific).
Two installations running the same model produce byte-identical \IR{1} graphs and byte-identical \IR{2} kernel graphs; their \IRstar{3} programs differ when their hardware differs; their binaries differ when their chips differ.
The cache hierarchy is therefore aligned with the IR hierarchy by construction, and federation safety follows from the upstream IRs' vendor-neutrality.

This concludes Part~III.
The next part (Numerical Portability and Determinism) elaborates the \struct{NumericalRecipe} contract, the recipe registry, and the cross-vendor numerics CI matrix that enforces the portability contract these IRs together establish.
The part after (Execution and Memory) elaborates the \struct{ExecutionPlan} layer that consumes \IRstar{3}-derived binaries and the static memory plan that backs them.
